\documentclass[12pt,a4paper]{article}
\usepackage[utf8]{inputenc}
\usepackage[T1]{fontenc}
\usepackage{lmodern}
\IfFileExists{brazil.ldf}{\usepackage[brazil]{babel}}{%
  \usepackage{babel}\babelprovide[main,import]{portuguese}}
\usepackage[margin=2.5cm,headheight=16.6pt]{geometry}
\usepackage{array,tabularx,booktabs}
\usepackage{amsmath}
\usepackage{enumitem}
\usepackage{fancyhdr}
\usepackage{xcolor}
\usepackage{titlesec}
\definecolor{azulProfessor}{HTML}{184C78}
\definecolor{cinzaTitulos}{HTML}{1E1E1E}
\definecolor{cinzaCabecalho}{gray}{0.5}
\titleformat{\section}{\Large\bfseries\color{azulProfessor}}{\thesection.}{0.6em}{}[\titlerule]
\titleformat{\subsection}{\large\bfseries\raggedright\color{cinzaTitulos}}{\thesubsection.}{0.6em}{}
\titleformat{\subsubsection}{\normalsize\bfseries\raggedright\color{cinzaTitulos}}{\thesubsubsection.}{0.6em}{}
\titlespacing*{\section}{0pt}{3.5ex plus 1ex minus .2ex}{2.3ex plus .2ex}
\titlespacing*{\subsection}{0pt}{3.25ex plus 1ex minus .2ex}{1.5ex plus .2ex}
\titlespacing*{\subsubsection}{0pt}{2.5ex plus 1ex minus .2ex}{1.2ex plus .2ex}
\setcounter{tocdepth}{2}
\usepackage{hyperref}
\hypersetup{colorlinks=true,linkcolor=azulProfessor,urlcolor=azulProfessor,citecolor=azulProfessor,linktoc=section,pdftitle={Laboratório Educativo de IA Generativa com GANs — Relatório Final — Sprint 1},pdfauthor={Allan Vitor Piran Lemes e Guilherme Rodrigues dos Santos}}
\urlstyle{same}
\newcolumntype{Y}{>{\raggedright\arraybackslash}X}
\setlength{\parindent}{1.5em}
\setlength{\parskip}{0pt}
\setlength{\emergencystretch}{2em}
\renewcommand{\arraystretch}{1.2}
\setlist[itemize]{leftmargin=*,itemsep=5pt,topsep=6pt}
\pagestyle{fancy}
\fancyhf{}
\fancyhead[L]{\small\color{cinzaCabecalho}Oficina de Projetos em Ciência de Dados e IA}
\fancyhead[R]{\small\color{cinzaCabecalho}Relatório Final -- UTFPR}
\fancyfoot[L]{\small\color{cinzaCabecalho}Documento de Ciência de Dados \& IA}
\fancyfoot[R]{\small Página \thepage}
\renewcommand{\headrulewidth}{0.4pt}
\renewcommand{\footrulewidth}{0.4pt}
\newcommand{\pendente}[1]{\textit{#1}}

\begin{document}
\renewcommand{\contentsname}{Sumário}
\renewcommand{\refname}{Referências}

\begin{titlepage}
\thispagestyle{fancy}
\centering
\vspace*{0.7cm}
{\Large\bfseries UNIVERSIDADE TECNOLÓGICA FEDERAL DO PARANÁ\par}
\vspace{0.15cm}
{\large DEPARTAMENTO ACADÊMICO DE INFORMÁTICA\par}
\vspace{0.1cm}
{\small CAMPUS PATO BRANCO\par}
\vspace{0.1cm}
{\small OFICINA DE PROJETOS EM CIÊNCIA DE DADOS E IA\par}
\vspace{1.6cm}
\noindent\rule{\textwidth}{1pt}\par
\vspace{0.2cm}
{\Huge\bfseries\color{azulProfessor}Laboratório Educativo de IA\par Generativa com GANs\par}
\vspace{0.35cm}
{\large\bfseries RELATÓRIO FINAL DE PROJETO\par}
\vspace{0.25cm}
{\small Atualização parcial: encerramento da Sprint 1\par}
\vspace{0.45cm}
\noindent\rule{\textwidth}{1pt}\par
\vspace{1.2cm}
{\renewcommand{\arraystretch}{1.05}
\begin{tabular}{@{}ll@{}}
\textbf{Equipe de Desenvolvimento:} & Allan Vitor Piran Lemes \\
 & Guilherme Rodrigues dos Santos \\
\textbf{Orientador / Prof.:} & Dr. Luiz Fernando Puttow Southier \\
\textbf{Data de Atualização:} & 9 de outubro de 2026 \\
\textbf{Versão do Documento:} & 1.2 \\
\end{tabular}\par}
\end{titlepage}
\setcounter{page}{2}
\tableofcontents
\clearpage

\section{Introdução}
Este relatório acompanha o desenvolvimento do Laboratório Educativo de IA Generativa com GANs, proposto para apresentar conceitos de inteligência artificial a estudantes do ensino médio por meio de uma apostila e de uma aplicação local. O projeto será desenvolvido por dois estudantes, com atividades de estudo, implementação e documentação distribuídas nas sprints previstas no Plano de Projeto entregue

A proposta consiste em utilizar uma Rede Adversarial Generativa (GAN) treinada com o Fashion-MNIST, uma base de imagens de roupas, calçados e acessórios. As imagens possuem resolução de $28\times28$ pixels, em tons de cinza~\cite{fashionmnist}. A aplicação educativa deverá mostrar exemplos da base, permitir a comparação de imagens geradas em diferentes momentos do treinamento e oferecer geração, visualização e download de imagens.

A utilidade do produto está em tornar o funcionamento da IA generativa mais acessível por meio de uma demonstração prática. A baixa resolução será apresentada como característica e limitação do modelo. O roteiro previsto terá duração inicial de 20 a 30 minutos e orientará a leitura e exploração autônoma por leitores iniciantes. O treinamento será realizado previamente pela equipe.

\textbf{Esta versão registra somente o encerramento da Sprint 1, de 02/10 a 09/10/2026.} As atividades relatadas compreendem a organização do trabalho, a preparação inicial dos dados e a definição dos leitores pretendidos. O desenvolvimento e o treinamento da GAN, a interface e a apostila permanecem como etapas posteriores. O relatório será atualizado ao final de cada sprint.


\section{Objetivos do Projeto}
\subsection{Objetivos}
O objetivo geral é desenvolver um laboratório educativo com uma GAN treinada no Fashion-MNIST, para apoiar a introdução à IA generativa no ensino médio por meio de materiais introdutórios, registrando o desenvolvimento e os aprendizados da equipe.

Os objetivos específicos do projeto são:
\begin{itemize}
\item Estudar o funcionamento da GAN, identificando o papel do gerador e do discriminador, e conhecer a base de dados utilizada.
\item Preparar o ambiente e adaptar uma implementação introdutória ao projeto.
\item Treinar o modelo e registrar a evolução visual das imagens geradas.
\item Construir uma interface simples para conhecer a base, comparar momentos do treinamento e gerar, visualizar e baixar imagens.
\item Documentar o código, as atividades, as dificuldades e os aprendizados de cada sprint.
\item Preparar material pedagógico e um roteiro de uso voltados ao ensino médio e a leitores iniciantes.
\end{itemize}

Esses objetivos abrangem o projeto completo. Na Sprint 1, o trabalho se concentrou na organização da equipe, no contato inicial com os dados e na caracterização da proposta educativa. A seleção de categorias de vestuário mencionada no planejamento corresponde a uma evolução posterior da geração básica e não integra as entregas desta sprint.

\subsection{Restrições Técnicas, Regulatórias e Éticas}
O desenvolvimento deve ser compatível com o tempo da disciplina, os recursos disponíveis e os conhecimentos adquiridos pelos integrantes durante o trabalho. Serão utilizados computadores pessoais e recursos gratuitos das ferramentas escolhidas. O uso de um ambiente de treinamento em nuvem será avaliado apenas se necessário e conforme sua disponibilidade.

O planejamento considera dedicação média de cinco horas semanais por integrante. As quatro sprints abrangem nove semanas, com mais duas semanas previstas para a consolidação do relatório e da apresentação. A data final de 18/12/2026 permanece sujeita à confirmação.

As fontes de dados e os exemplos de código efetivamente utilizados deverão ser identificados. As imagens geradas serão apresentadas como resultados de IA, incluindo suas limitações. A avaliação prevista é descritiva: funcionamento da aplicação, evolução visual e revisão dos materiais pelos autores. Não foi estabelecida uma porcentagem de acerto para a GAN.

As 10.000 imagens do conjunto oficial de teste do Fashion-MNIST devem permanecer fora do treinamento. A divisão adicional realizada na Sprint 1 utiliza apenas o conjunto oficial de treinamento, conforme detalhado na Seção~\ref{sec:dados}.

\subsection{Avaliação e Acurácia das Estimativas (Planejado vs. Executado)}
O plano adotou uma estimativa baseada na capacidade da equipe. Foram previstas 110 horas-pessoa para o projeto: 90 horas nas quatro sprints e 20 horas para a etapa final. Esse total representa uma previsão, e não o esforço já realizado. Para esta atualização, a comparação se limita à Sprint 1, cuja capacidade prevista é de 10 horas-pessoa, incluindo o planejamento inicial.

\begin{center}
\small
\begin{tabularx}{\textwidth}{@{}YrrY@{}}
\toprule
\textbf{Integrante} & \textbf{Horas previstas} & \textbf{Custo ilustrativo} & \textbf{Horas realizadas} \\
\midrule
Allan Vitor Piran Lemes & 5 & R\$ 100,00 &  3 \\
Guilherme Rodrigues dos Santos & 5 & R\$ 100,00 & 2 \\
\midrule
\textbf{Total da Sprint 1} & \textbf{10} & \textbf{R\$ 200,00} & \textbf{5} \\
\bottomrule
\end{tabularx}
\end{center}

O custo ilustrativo segue o valor de R\$ 20,00 por hora definido no plano. Assim, o projeto completo possui estimativa de R\$ 2.200,00, dos quais R\$ 200,00 correspondem à Sprint 1. Esses valores atendem ao exercício de planejamento e não representam pagamento aos estudantes. O desembolso direto previsto é de R\$ 0,00, com uso dos recursos já disponíveis; esta versão não apresenta um levantamento de despesas efetivamente realizadas.

A variação do esforço será calculada após a consolidação das horas individuais:
\[
\text{Desvio (\%)} =
\frac{\text{Horas realizadas} - \text{Horas previstas}}
     {\text{Horas previstas}} \times 100.
\]

Não foi calculado um desvio nesta atualização porque as horas realizadas não foram informadas. O período de referência encerra-se em 09/10/2026; a avaliação de cumprimento das tarefas está na Seção~\ref{sec:comparacao}. Eventuais diferenças de esforço deverão ser explicadas com base nas atividades e dificuldades efetivamente registradas. O tempo em que o computador executa tarefas sozinho não deve ser contabilizado como dedicação ativa da equipe.

\section{Gestão de Riscos do Projeto}
Os quatro riscos abaixo correspondem ao planejamento entregue. Nesta versão, sua análise é parcial e se limita à Sprint 1. A ausência de um registro de ocorrência não será interpretada como confirmação de que o risco não ocorreu.

\subsection*{R01 — Dificuldade para compreender ou executar a GAN}
\textbf{Previsão:} probabilidade alta e impacto alto. A prevenção consiste em começar com um exemplo introdutório, estudar em conjunto e testar cada etapa. Caso necessário, a equipe buscará apoio do professor e adotará uma implementação menor, preservando a geração básica.

\textbf{Situação na Sprint 1:} o registro contempla a preparação dos dados, sem resultados de implementação ou treinamento da GAN. A ocorrência desse risco e a eficácia das ações relacionadas ao modelo ainda não podem ser avaliadas.

\subsection*{R02 — Treinamento demorado ou interrompido}
\textbf{Previsão:} probabilidade média e impacto médio. Estão previstos um teste curto antes de execuções maiores e o salvamento do progresso. A contingência é retomar a versão salva ou avaliar outro ambiente gratuito disponível, reduzindo o tempo de treinamento se necessário.

\textbf{Situação na Sprint 1:} não há treinamento registrado nesta etapa. Ainda não existem informações sobre duração, interrupções ou necessidade de mudança de ambiente.

\subsection*{R03 — Imagens pouco reconhecíveis}
\textbf{Previsão:} probabilidade alta e impacto médio. A prevenção é conferir os dados e acompanhar as imagens desde os primeiros testes. Caso o risco ocorra, serão revisados o exemplo utilizado e pequenos ajustes, registrando as limitações e preservando tempo para concluir a aplicação.

\textbf{Situação na Sprint 1:} foi relatada a exportação de um painel com imagens da base. Esse painel apoia o contato inicial com os dados, mas não permite avaliar a qualidade de imagens geradas, pois não há resultados da GAN nesta atualização.

\subsection*{R04 — Falta de tempo para concluir as atividades}
\textbf{Previsão:} probabilidade média e impacto alto. A prevenção é acompanhar as tarefas no Jira e comparar o trabalho restante com a capacidade da sprint. A contingência é adiar melhorias opcionais e priorizar a demonstração educativa, as funções básicas e os registros das entregas.

\textbf{Situação na Sprint 1:} o quadro inicial do Jira e os prazos das quatro sprints foram organizados. Essa ação fornece uma base para o acompanhamento, mas ainda não há horas consolidadas ou histórico suficiente para avaliar sua eficácia na prevenção de atrasos. Não foi apresentado um registro específico de ocorrência desse risco.

\section{Execução por Sprints e Alocação de Recursos}
O desenvolvimento está organizado em quatro sprints e uma etapa de consolidação final. A equipe é formada por Allan Vitor Piran Lemes e Guilherme Rodrigues dos Santos. Os percentuais associados às entregas representam pesos de avaliação da disciplina, e não percentuais de conclusão ou de esforço.

\subsection{Primeira Sprint: Organização, Jira e Contato com os Dados}

O objetivo desta sprint foi estabelecer a organização inicial do trabalho, preparar o acesso à base Fashion-MNIST e relacionar a proposta técnica à linguagem acessível ao ensino médio. As atividades descritas a seguir correspondem ao registro de execução da equipe.

\subsubsection{Atividades realizadas}
Foi configurado o repositório Git no GitHub, incluindo o arquivo \path{.gitignore}, para organizar o código e controlar quais arquivos serão acompanhados no versionamento. Também foi estruturado o quadro inicial do Jira, com as quatro sprints e os respectivos prazos. Essa organização servirá de referência para distribuir e acompanhar as próximas tarefas.

Na preparação dos dados, foi desenvolvido o script Python \path{preparacao_dados.py} para baixar o Fashion-MNIST. O registro da sprint informa a separação de imagens para treinamento e validação e a exportação de um painel inicial de amostras da base no diretório \path{artefatos/}. As características dessa organização são apresentadas a seguir.

Na dimensão educativa, foi redigido o objetivo do material introdutório com linguagem acessível ao ensino médio.

\subsubsection{Organização inicial dos dados}
\label{sec:dados}
O Fashion-MNIST possui 60.000 imagens no conjunto oficial de treinamento e 10.000 no conjunto oficial de teste, distribuídas em dez categorias~\cite{fashionmnist}. Nesta sprint, o conjunto de treinamento foi subdividido em 50.000 imagens para treinamento e 10.000 para validação, com divisão estratificada e semente 42. A estratificação busca manter a proporção das categorias nos dois subconjuntos; a semente permite repetir a mesma divisão.

\begin{center}
\small
\begin{tabularx}{\textwidth}{@{}YrY@{}}
\toprule
\textbf{Conjunto} & \textbf{Quantidade} & \textbf{Origem e destinação} \\
\midrule
Treinamento & 50.000 & Parte do treinamento oficial, destinada ao aprendizado da GAN. \\
Validação & 10.000 & Parte reservada do treinamento oficial para verificações posteriores. \\
Teste oficial & 10.000 & Conjunto original separado, mantido fora do treinamento. \\
\bottomrule
\end{tabularx}
\end{center}

Essa subdivisão é um detalhamento da preparação dos dados que não estava especificado no plano. Ela mantém a separação do teste oficial. Sua existência não significa que o modelo já foi treinado ou avaliado; o procedimento de uso da validação será definido durante a implementação. O painel exportado contém exemplos da base, e não imagens produzidas pela GAN.

\subsubsection{Comparação entre o planejado e o realizado}
\label{sec:comparacao}
\begin{center}
\small
\begin{tabularx}{\textwidth}{@{}>{\raggedright\arraybackslash}p{0.29\textwidth}Y@{}}
\toprule
\textbf{Previsão da Sprint 1} & \textbf{Registro no encerramento} \\
\midrule
Organizar o repositório e preparar o ambiente & Repositório configurado, com \path{.gitignore}, e preparação dos dados realizada em Python nos computadores da equipe. \\
Organizar o Jira e registrar os prazos & Quadro inicial configurado com quatro sprints e prazos. O registro disponível não detalha os cartões e seus critérios de conclusão. \\
Definir tarefas e responsáveis & Divisão principal entre os integrantes registrada. O detalhamento dos cartões deverá ser vinculado como evidência. \\
Baixar a base e visualizar exemplos & Script de preparação desenvolvido, divisão dos dados realizada e painel de amostras exportado. \\
Descrever a proposta educativa & Objetivo educativo e leitores pretendidos definidos. \\
Registrar atividades, dúvidas e aprendizados & Atividades e contribuições reunidas neste relatório. Faltam as horas, as reflexões individuais e a inclusão das evidências. \\
\bottomrule
\end{tabularx}
\end{center}

\subsubsection{Recursos utilizados e divisão das tarefas}
Foram utilizados os computadores pessoais dos integrantes. As ferramentas informadas nesta sprint foram Python 3.10 ou superior, PyTorch, Torchvision, Scikit-learn, Git/GitHub, Jira e Overleaf. Não há uso de infraestrutura de treinamento em nuvem registrado nesta etapa.

\begin{center}
\small
\begin{tabularx}{\textwidth}{@{}>{\raggedright\arraybackslash}p{0.28\textwidth}Y@{}}
\toprule
\textbf{Integrante} & \textbf{Contribuição na Sprint 1} \\
\midrule
Allan Vitor Piran Lemes & Organização do repositório, desenvolvimento da preparação dos dados e caracterização dos leitores da apostila. \\
Guilherme Rodrigues dos Santos & Estruturação do quadro inicial do Jira e planejamento das tarefas das sprints. \\
\bottomrule
\end{tabularx}
\end{center}

A distribuição das atividades está de acordo com a divisão principal prevista para a primeira sprint no planejamento entregue. A organização adotada permite que cada integrante conduza suas tarefas e acompanhe o trabalho do colega. As horas efetivamente dedicadas ainda devem ser consolidadas na Seção 2.3.

\subsubsection{Aprendizados, dificuldades e pendências}
As atividades desta etapa envolveram três temas de aprendizagem: a organização do código em um repositório, a divisão do trabalho em tarefas e prazos e a preparação de uma base de imagens para uso posterior. A separação dos conjuntos também introduziu a distinção entre dados destinados ao treinamento e dados reservados para verificação. No contexto educativo, foi necessário definir uma linguagem acessível para os materiais.

O registro disponível descreve as atividades, mas não apresenta uma reflexão individual sobre os conceitos compreendidos, as dúvidas ou os problemas enfrentados.

Permanecem pendentes de documentação as horas individuais, os links de acesso ao código e ao quadro e a inclusão do painel de amostras.

\subsubsection{Evidências da Sprint 1}
Os artefatos abaixo foram mencionados no registro de execução. Seus links e imagens devem ser incorporados para permitir a consulta do trabalho realizado.
\begin{itemize}
\item \textbf{Repositório:} \pendente{inserir o endereço do GitHub e a referência da versão correspondente ao encerramento da sprint.}
\item \textbf{Preparação dos dados:} script \path{preparacao_dados.py}. \pendente{Inserir o link para o arquivo e, se disponível, o registro de sua execução.}
\item \textbf{Painel de amostras:} imagem exportada no diretório \path{artefatos/}. \pendente{Inserir o painel e seu nome completo.}
\item \textbf{Jira:} \pendente{inserir uma captura ou exportação do quadro com as sprints, os prazos, os responsáveis e a situação das tarefas no encerramento.}
\end{itemize}

\subsubsection{Encaminhamento para a próxima sprint}
A continuidade prevista é estudar uma implementação introdutória de GAN, adaptar o código e executar um primeiro treinamento. Esses encaminhamentos seguem o planejamento entregue; não representam atividades concluídas na Sprint 1.

\section{Resultados Obtidos e Verificação dos Materiais}
\subsection{Resultados Parciais da Sprint 1}
Os resultados registrados nesta etapa são o repositório organizado, o quadro inicial do Jira, o script de preparação do Fashion-MNIST, a divisão dos dados e o painel de exemplos da base. Também foram descritos o objetivo educativo e os leitores pretendidos. Essas entregas estabelecem as condições iniciais para o desenvolvimento do modelo e para o acompanhamento do projeto.

O relato da equipe indica avanço nas atividades de organização e preparação previstas. Sua documentação ainda precisa ser complementada com evidências, horas realizadas e reflexões individuais.

\begin{thebibliography}{9}
\addcontentsline{toc}{section}{Referências}

\bibitem{fashionmnist}
XIAO, Han; RASUL, Kashif; VOLLGRAF, Roland.
\textbf{Fashion-MNIST: a Novel Image Dataset for Benchmarking Machine Learning Algorithms}.
2017. Repositório do projeto disponível em:
\url{https://github.com/zalandoresearch/fashion-mnist}.
\end{thebibliography}

\end{document}
